\documentclass[letterpaper]{article} % DO NOT CHANGE THIS
% Supplement, AAAI-27 Author Kit.
\usepackage{aaai2027}  % DO NOT CHANGE THIS
\usepackage[hyphens]{url}  % DO NOT CHANGE THIS
\usepackage{graphicx} % DO NOT CHANGE THIS
\urlstyle{rm} % DO NOT CHANGE THIS
\def\UrlFont{\rm}  % DO NOT CHANGE THIS
\usepackage{natbib}  % DO NOT CHANGE THIS AND DO NOT ADD ANY OPTIONS TO IT
\usepackage{caption} % DO NOT CHANGE THIS AND DO NOT ADD ANY OPTIONS TO IT
\frenchspacing  % DO NOT CHANGE THIS
\usepackage{booktabs}
\usepackage{amsmath}
\usepackage{array}
\newcommand{\lbl}{\textsc{Label}}
\pdfinfo{
/TemplateVersion (2027.1)
}
\setcounter{secnumdepth}{2}
\renewcommand{\thesection}{\Alph{section}}

\title{Supplementary Material for\\\lbl{}: Literature Asymmetry between Efficacy and Limits\\in Drug Knowledge and Language Models}
\author{S M Rayeed, Yao Ma}
\affiliations{Rensselaer Polytechnic Institute, Troy, NY, USA}

\begin{document}
\maketitle

\section{Measuring Whether Text Names a Fact}

\paragraph{Names.}
Each disease has three sets of names. The exact name is its MONDO label
\citep{vasilevsky2022mondo}. Alternative names are the titles of Wikipedia
redirects to the disease's article, which include synonyms, lay terms and
abbreviations. Family names are the labels of the disease's MONDO ancestors up
to two levels above it, taken from a directed child-to-parent table of 17,801
edges over 11,788 diseases. The sets are cumulative: a text that names the
disease at a finer level is also counted at every coarser one, and each edge is
assigned the finest level it reaches.

\paragraph{Matching rule.}
A phrase matches when it occurs on word boundaries, ignoring case, with an
optional trailing ``s'' to allow plurals. Multi-word phrases match across any
run of whitespace. Phrases shorter than four characters are never matched, so
short abbreviations do not produce coincidental hits. Under this rule
``asthma'' does not match ``asthmatic,'' and ``colitis'' does not match inside
a longer word.

\paragraph{Floor.}
For every edge, three diseases that the drug has no edge to in that relation
are drawn from the same relation's targets, and the same matcher is run on the
drug's text for them. The share of these unrelated diseases that match is the
floor. It estimates how often a disease name appears in a drug's text by
coincidence rather than as a statement of the relation. Table~\ref{tab:floor}
gives the exact-name floor on the probe sample.

\begin{table}[h]
\centering
\small
\begin{tabular}{lcc}
\toprule
Relation & Real edge & Floor \\
\midrule
Indication & 17.17 & 1.22 \\
Off-label use & 10.50 & 0.72 \\
Contraindication & 5.50 & 0.78 \\
\bottomrule
\end{tabular}
\caption{Exact-name match rate (\%) in Wikipedia for real edges and for three
unrelated diseases per edge, 600 edges per relation.}
\label{tab:floor}
\end{table}

\paragraph{Contribution of each level.}
On the same draw in Wikipedia, averaged over the three relations, adding
alternative names to the exact name raises the match rate by 1.39 points.
Adding the disease's MONDO parent raises it by a further 5.50 points, and
adding ancestors up to two levels by 8.43 points. Each increase is the
difference of floor-corrected rates, with the floor recomputed at every
level, because wider name sets also match unrelated diseases more often. Only the exact and
alternative names count as naming the fact in the main paper.

\section{Corpora}

\paragraph{Wikipedia.}
For each drug we fetch the current English Wikipedia article by title and
cache its source text. The text measurement runs on the article as fetched.

\paragraph{PubMed.}
Abstracts are retrieved through the NCBI E-utilities \citep{sayers2009eutils}
with the query \texttt{"<drug>"[Title/Abstract] AND (<fragment>)}, sorted by
relevance. Each drug receives a budget of 40 abstracts, split evenly between
two fragments, one framed for therapeutic use and one for adverse effects and
contraindications, so that each relation receives 20. The union of returned
records is fetched once and cached. A request that fails is not cached, so a
failure is never recorded as an empty result. The first version of the corpus
used a single relevance-ranked query of 40 abstracts per drug; the comparison
of the two versions is reported in the main paper.

\section{Model Probe}

\paragraph{Question format.}
Each edge is scored as a completion. The prompt is a relation-specific stem
that names the drug and stops before the disease (Table~\ref{tab:stems}); each
option is appended as a continuation, and its score is the mean token
log-probability of the continuation. The option with the highest score is the
model's answer. No chat template, passage or retrieval is used.

\begin{table}[h]
\centering
\small
\begin{tabular}{>{\raggedright\arraybackslash}p{1.75cm}>{\raggedright\arraybackslash}p{5.35cm}}
\toprule
Relation & Stem \\
\midrule
Indication & \textit{X is a drug used to treat the disease called} \\
Contraind. & \textit{X is a drug that must not be given to patients who have} \\
Off-label use & \textit{X is a drug sometimes used off-label to treat} \\
Treats (Hetionet) & \textit{X is a drug used to treat the disease called} \\
Palliates (Hetionet) & \textit{X is a drug used to relieve the symptoms of} \\
\bottomrule
\end{tabular}
\caption{Prompt stems; X is the drug's name.}
\label{tab:stems}
\end{table}

\paragraph{Options.}
Each question has five options: the true disease and four distractors drawn
at random from the same relation's targets in the sample. Distractors exclude
every disease the drug is paired with in that relation within the sample.
Partners outside the sample cannot be excluded; with about 600 targets per
relation and a handful of partners per drug, the chance that a distractor is
an unrecorded true partner is on the order of one percent, and it applies
equally to all relations. The order of options is shuffled with a fixed seed.

\paragraph{Models.}
Qwen2.5-0.5B-Instruct and Qwen2.5-7B-Instruct \citep{qwen2025qwen25} run in
32-bit floating point on CPU with Hugging Face Transformers.

\section{Prompt Stems and Declines}

\begin{table}[h]
\centering
\footnotesize
\setlength{\tabcolsep}{3pt}
\begin{tabular}{p{2.15cm}p{4.95cm}}
\toprule
Relation & Prompt stem (X is the drug's name) \\
\midrule
Indication & \textit{X is a drug used to treat the disease called} \\
Contraindication & \textit{X is a drug that must not be given to patients who have} \\
Off-label use & \textit{X is a drug sometimes used off-label to treat} \\
Treats & \textit{X is a drug used to treat the disease called} \\
Palliates & \textit{X is a drug used to relieve the symptoms of} \\
\bottomrule
\end{tabular}
\caption{Prompt stems for the closed-book probe. Each option is scored as the
continuation of the stem; the last two rows are for the Hetionet relations.}
\label{tab:stems}
\end{table}

\begin{table}[h]
\centering
\footnotesize
\setlength{\tabcolsep}{3.5pt}
\begin{tabular}{lcccc}
\toprule
 & \multicolumn{2}{c}{Indication} & \multicolumn{2}{c}{Contraindication} \\
\cmidrule(lr){2-3}\cmidrule(lr){4-5}
Page & closed & page & closed & page \\
\midrule
Names the disease & 6.0 & 5.8 & 36.2 & 41.6 \\
Names only its family & 9.1 & 26.0 & 42.1 & 54.4 \\
Silent & 29.6 & 55.6 & 43.3 & 76.2 \\
\bottomrule
\end{tabular}
\caption{Share of questions answered ``Not enough information'' (\%), closed
book and with the drug's page, averaged over Qwen2.5 7B, OLMo 2 7B, and
Granite 3.1 8B.}
\label{tab:declines}
\end{table}

\section{Negative Control on Hetionet}

\begin{table}[h]
\centering
\footnotesize
\setlength{\tabcolsep}{3pt}
\begin{tabular}{lcc}
\toprule
 & \emph{treats} & \emph{palliates} \\
\midrule
Named in Wikipedia (\%) & 57.33 (n=600) & 56.10 (n=385) \\
Odds ratio, pooled & \multicolumn{2}{c}{1.05, $p = 0.70$} \\
Odds ratio, within drug (51) & \multicolumn{2}{c}{2.48 [1.29, 4.75]} \\
Relation term, degree held fixed & \multicolumn{2}{c}{$p = 0.16$} \\
\midrule
Gap by edges per drug: 1--2 & \multicolumn{2}{c}{$+8.16$ ($p = 0.06$)} \\
\phantom{Gap by edges per drug:} 3--5 & \multicolumn{2}{c}{$+0.93$ ($p = 0.89$)} \\
\phantom{Gap by edges per drug:} 6--12 & \multicolumn{2}{c}{$-11.38$ ($p = 0.32$)} \\
\bottomrule
\end{tabular}
\caption{Negative control on Hetionet. Two benefit relations, same corpus and
matcher. No gap survives the degree control, and the strata change sign.}
\label{tab:hetionet}
\end{table}

\section{Harder-Distractor Check}

\paragraph{Questions.}
The check reuses the question set of the companion abstention study. Each
question asks ``Which disease is X used to treat?'' (indication) or ``In
patients with which disease must X not be given?'' (contraindication) and
lists five diseases and the option ``Not enough information to tell,''
lettered A to F. The position of each option, including the decline option,
is random.

\paragraph{Distractors.}
Wrong options are taken first from the true disease's MONDO siblings, then
from its cousins (other children of its grandparent), and only then from
random diseases of the same relation. They never include the true disease's
own ancestors or descendants, or any disease the drug is related to by the
same relation. In 481 of 856 questions all four distractors come from siblings
or cousins.

\paragraph{Scoring.}
Each prompt is formatted with the model's chat template, and the answer is the
option letter with the highest next-token probability. Both spellings of each
letter, with and without a leading space, are scored and the better one kept;
probabilities are renormalized over the six letters. Prompts are batched with
left padding, and position indices are rebuilt from the attention mask so that
batched and single scores agree.

\paragraph{Models.}
Qwen2.5 0.5B and 7B, OLMo~2 7B \citep{olmo2024olmo2} and Granite~3.1 8B
\citep{granite2024}, all instruction-tuned, in 32-bit floating point on CPU.

\section{Statistical Procedures}

\paragraph{Within-drug odds ratio.}
For drug $d$, let $a_d$ and $b_d$ be its named and unnamed indication edges,
$c_d$ and $e_d$ its named and unnamed contraindication edges, and $n_d$ their
total. The Mantel--Haenszel estimate \citep{mantel1959} is
\begin{equation}
\mathrm{OR}_{\mathrm{MH}} = \frac{\sum_d a_d e_d / n_d}{\sum_d b_d c_d / n_d},
\end{equation}
with a 95\% confidence interval $\exp(\log \mathrm{OR}_{\mathrm{MH}} \pm
1.96/\sqrt{\sum_d V_d})$, where $V_d = n_{d,1} n_{d,2} m_{d,1} m_{d,2} / (n_d^2
(n_d - 1))$ is the hypergeometric variance of drug $d$'s table with row totals
$n_{d,\cdot}$ and column totals $m_{d,\cdot}$, and a continuity-corrected
Cochran--Mantel--Haenszel test of no association.

\paragraph{Edge-count control.}
A logistic model regresses whether a fact is named on an indicator for the
relation and on $\log$ of the drug's number of edges in that relation in the
full graph. The relation coefficient tests the gap with edge count held fixed.

\paragraph{Wilson intervals.}
For $k$ successes in $n$ trials, $\hat p = k/n$, and $z = 1.96$, the interval
is
\begin{equation}
\frac{\hat p + \frac{z^2}{2n} \pm z\sqrt{\frac{\hat p(1-\hat p)}{n} +
\frac{z^2}{4n^2}}}{1 + \frac{z^2}{n}}.
\end{equation}

\paragraph{Paired gain from scale.}
For each edge $i$, the gain is $g_i = y_i^{7B} - y_i^{0.5B}$, where $y = 1$
for a correct answer. The statistic is the mean gain on indications minus the
mean gain on contraindications. Its interval comes from 5,000 bootstrap
resamples in which drugs, not edges, are drawn with replacement, so that edges
of the same drug move together. The $p$-value uses the normal approximation
with the bootstrap standard error.

\paragraph{Gap within naming groups.}
In the distractor check, the indication-minus-contraindication difference is
computed separately within questions whose text names the answer, names only
its family, or names nothing. The group differences $\delta_g$ are combined as
$\sum_g w_g \delta_g / \sum_g w_g$ with $w_g = n_{g,1} n_{g,2} / (n_{g,1} +
n_{g,2})$, and the standard error follows from the binomial variances of the
group proportions.

\section{Additional Results}

\begin{table}[h]
\centering
\small
\setlength{\tabcolsep}{2pt}
\begin{tabular}{lccc}
\toprule
Relation & 0.5B [95\% CI] & 7B [95\% CI] & Gain \\
\midrule
Indication & 31.50 [27.91, 35.32] & 50.83 [46.84, 54.82] & +19.33 \\
Off-label & 28.50 [25.03, 32.24] & 35.00 [31.29, 38.90] & +6.50 \\
Contraind. & 22.17 [19.03, 25.66] & 23.83 [20.60, 27.40] & +1.67 \\
\bottomrule
\end{tabular}
\caption{Probe accuracy (\%) and paired gain from 0.5B to 7B (points), 600
edges per relation. Indication gain minus contraindication gain: +17.67 [12.68,
22.70], $p = 4.8\times10^{-12}$.}
\label{tab:probe_full}
\end{table}

\begin{table}[h]
\centering
\small
\begin{tabular}{lccc}
\toprule
Edges per drug & 3--5 & 6--12 & 13+ \\
\midrule
Wikipedia & +13.94 & +6.04 & +1.64 \\
Wikipedia + PubMed & +11.81 & +6.90 & +2.23 \\
\bottomrule
\end{tabular}
\caption{Indication minus contraindication coverage (points) within strata of
the drug's number of edges in the relation. The 13+ stratum is not
distinguishable from zero in either corpus.}
\label{tab:degree}
\end{table}

\begin{table}[h]
\centering
\small
\setlength{\tabcolsep}{4pt}
\begin{tabular}{lccccc}
\toprule
Model & Ind. & Contra. & Gap & Within & $p$ (within) \\
\midrule
Qwen2.5 0.5B & 36.0 & 28.1 & +7.9 & +8.2 & $8.9\times10^{-3}$ \\
Qwen2.5 7B & 56.0 & 25.1 & +30.9 & +31.3 & $<10^{-15}$ \\
OLMo 2 7B & 56.0 & 33.0 & +23.0 & +23.6 & $2.3\times10^{-13}$ \\
Granite 3.1 8B & 58.7 & 12.8 & +45.9 & +46.3 & $<10^{-15}$ \\
\bottomrule
\end{tabular}
\caption{Harder-distractor check, closed book (\%). 450 indication and 406
contraindication questions. Raw-gap $p$-values are $0.013$ for Qwen2.5 0.5B,
$1.5\times10^{-11}$ for OLMo~2 7B, and below $10^{-15}$ for Qwen2.5 7B and
Granite~3.1 8B.}
\label{tab:distractor_full}
\end{table}

\paragraph{Control graph.}
On Hetionet, Wikipedia names the disease for 57.33\% of 600 sampled
\emph{treats} edges and 56.10\% of 385 measurable \emph{palliates} edges, an
odds ratio of 1.05 ($p = 0.70$). Within the 51 drugs that carry both
relations, the Mantel--Haenszel odds ratio is 2.48 (95\% CI 1.29--4.75). With
the drug's edge count held fixed, the relation term is not significant
($p = 0.16$), and the edge-count strata show no consistent direction (+8.2
points for 1--2 edges, +0.9 for 3--5, $-11.4$ for 6--12). In PrimeKG, the
relation term remains significant under the same control ($p = 2.4\times10^{-4}$
in Wikipedia). Hetionet uses the Disease Ontology rather than MONDO for its
hierarchy.

\section{Reproducibility}

\paragraph{Pipeline.}
The analysis runs in five stages, each a standalone script that reads the
previous stage's files: building the edge samples from PrimeKG, Hetionet and
MONDO; fetching and caching Wikipedia articles and PubMed abstracts; measuring
the naming level of every edge against its floor; scoring the model probe; and
computing the within-drug, paired and distractor statistics. A single
script runs all five stages from the public downloads, submits the model jobs
to the cluster scheduler, and writes every table. The statistics use the
Python standard library only. Random draws use a fixed seed of 0.

\paragraph{Compute.}
All model scoring runs on CPU nodes of a university supercomputer, with
PyTorch 2.14 and Transformers 5.20 in 32-bit floating point. No GPU is
required.

\paragraph{Data.}
PrimeKG, Hetionet, MONDO, Wikipedia and PubMed are public. We release the
sampled edge lists, the per-edge naming levels and the per-edge model scores,
together with the scripts, so that every table in the paper and this
supplement can be regenerated from the released files.

\bibliography{refs}

\end{document}
